\chapter{Solver, Fehlermessung und Wärmeindikatoren}\label{chap:solvers}
\section{Kongruente Skalierung und direkte Referenz}
Nach Rand- und Elektrodenreduktion löst die Referenz das freie eindeutige System
$\mathbf A\mathbf x=\mathbf b$. Bei positiver Diagonalmatrix $\mathbf S$ gilt
\begin{equation}\label{eq:scale}
(\mathbf S\mathbf A\mathbf S)\mathbf y=\mathbf S\mathbf b,
\qquad\mathbf x=\mathbf S\mathbf y.
\end{equation}
Diese Kongruenz erhält Symmetrie und Trägheit der Matrix; sie ändert weder
Physik noch Eindeutigkeit. Einzelne positive Feldsysteme verwenden die
Kehrwurzel der Diagonalen. Piezo verwendet feldweise mediane positive
mechanische/dielektrische Diagonalen $a_u$ und $a_\phi$ mit
$S_u=a_u^{-1/2}$, $S_\phi=a_\phi^{-1/2}$.
N/m und F dürfen ohne diese Einheitenbehandlung nicht pauschal verglichen werden.
\index{Skalierung}\index{Sparse LU}

SciPy Sparse LU ist die Produktionsreferenz über \texttt{splu}; die
SciPy-Dokumentation \cite{scipy-lu} erläutert den direkten CSC-/LU-Einstieg.
Faktoren besitzen Kopien und können für mehrere rechte Seiten wiederverwendet
werden. Komplexe rechte Seiten zu reellen Faktoren werden mit getrennten
Real-/Imaginärsolves behandelt. Der Standardpfad hat keinen stillen Fallback.
Pivotgrenzen diagnostizieren numerische Singularität, sind aber kein
vollständiger Konditionszahlnachweis. CG ist nur für wirklich positiv definite
Einzelfelder mathematisch zulässig und kein auswählbarer Ersatz des
monolithischen Piezo-Solvers. MINRES ist ohne geeigneten symmetrischen
Präconditioner nicht als unterstützt zu beschreiben.

\section{PETSc direkt und FieldSplit}
Der optionale Direktadapter verwendet in einem Prozess \texttt{COMM\_SELF},
SeqAIJ, PREONLY/LU. Core-Import benötigt kein PETSc; die Konstruktion prüft die
native Laufzeit. Model/TOML verwenden SciPy; direkte Python-Einstiege wählen
den Backend explizit. Ein realer PETSc-Build bildet komplexe Systeme durch
\begin{equation}\label{eq:complex-embed}
\begin{pmatrix}\operatorname{Re}\mathbf A&-\operatorname{Im}\mathbf A\\
\operatorname{Im}\mathbf A&\operatorname{Re}\mathbf A\end{pmatrix}
\begin{pmatrix}\operatorname{Re}\mathbf x\\\operatorname{Im}\mathbf x\end{pmatrix}
=\begin{pmatrix}\operatorname{Re}\mathbf b\\\operatorname{Im}\mathbf b\end{pmatrix}
\end{equation}
ab; komplexe PETSc-Builds arbeiten direkt. Beide erhalten sämtliche
komplexen Matrix- und Lastanteile.
\index{PETSc}\index{FieldSplit}

Für zwei nichtleere Felder besitzt
$\mathbf A=\left[\begin{smallmatrix}\mathbf A_{00}&\mathbf A_{01}\\
\mathbf A_{10}&\mathbf A_{11}\end{smallmatrix}\right]$ die Faktorisierung
\begin{equation}\label{eq:schur-factor}
\mathbf A=
\begin{pmatrix}\mathbf I&0\\\mathbf A_{10}\mathbf A_{00}^{-1}&\mathbf I\end{pmatrix}
\begin{pmatrix}\mathbf A_{00}&0\\0&\mathbf S_c\end{pmatrix}
\begin{pmatrix}\mathbf I&\mathbf A_{00}^{-1}\mathbf A_{01}\\0&\mathbf I\end{pmatrix},
\end{equation}
mit dem Schur-Komplement
\begin{equation}
\mathbf S_c=\mathbf A_{11}-\mathbf A_{10}\mathbf A_{00}^{-1}\mathbf A_{01}.
\end{equation}
Der implementierte FieldSplit verwendet volle Schur-Blockfaktorisierung,
LU für $\mathbf A_{00}$ und $\mathbf A_{11}^{-1}$ als Approximation für
$\mathbf S_c^{-1}$. Der Schur-Operator wird nicht dicht assembliert.
Die offizielle FieldSplit-Quelle \cite{petsc-fieldsplit} beschreibt die
Blockstruktur; die konkrete A11-Wahl ist der Projektvertrag aus Doc/39.
Singuläre Diagonalblöcke werden nicht durch Verschiebungen repariert.

Der äußere Solver ist FGMRES mit Rechtspräkonditionierung. Die approximative
Blockfaktorisierung ist nicht notwendigerweise symmetrisch oder positiv;
deshalb passen CG/MINRES hier nicht. Die FGMRES-Quelle \cite{petsc-fgmres}
dokumentiert die unterstützte Präkonditionierungsseite. Ein nichtleeres Feld
verwendet FGMRES mit LU, ein leeres System null Iterationen, mehr als zwei
Felder sind Fehler. Feldpartitionen beziehen sich auf das kondensierte freie
System; Real-/Imaginärteile eines Feldes werden bei Einbettung nicht als
unterschiedliche Physikfelder behandelt. Grenzen, Neustarts und Toleranzen
sind explizite Parameter. Nichtkonvergenz wird als Fehler weitergegeben.

Die unabhängigen kleinen Paritätstests und der vorhandene Vier-Netze-Wärmefall
prüfen reale/komplexe und direkte/iterative Pfade. Sie sind kein MPI-, AMG-
oder allgemeiner Großmodell-Skalierungsnachweis. Getestete native Umgebung:
PETSc 3.19.6 in der separaten CI; aktuelle Online-Dokumentation ist eine
Methodenquelle, keine Behauptung eines Versionsupdates der Toolbox.

\section{Physikalische Residuen und Rundungsbudget}
Die unskalierte Originalgleichung liefert $\mathbf r=\mathbf A\mathbf x-\mathbf b$.
Für einen Backend-Feldindexsatz $I$ lautet der Rückwärtsfehler
\begin{equation}\label{eq:backward-error}
\eta_I=\frac{\|\mathbf r_I\|_2}
{\|(|\mathbf A||\mathbf x|+|\mathbf b|)_I\|_2}.
\end{equation}
Bei Nenner null ist das Verhältnis null nur für Residuum null, sonst unendlich.
Normen werden intern gegen gemeinsamen Über-/Unterlauf skaliert.
Backend und Physikadapter verlangen geeignete freie bzw. kondensierte
Feldresiduen; die Backend-Grenze ist $10^{-9}$.
Einzelfeld-Physikberichte besitzen ihre spezifischen dokumentierten
Nenner-/Freiselektionen; sie sind nicht immer numerisch identisch mit dem
Backendbericht. KSP-Erfolg prüft zunächst das skalierte System, danach folgt
diese Originalkontrolle. Mechanische Trägheit gehört im dynamischen Residuum
zur Originalgleichung. Bilanztoleranzen kombinieren physikalische Lastskalen
mit einem Rundungsterm aus absoluten, ungekürzten Matrixprodukten.
Sie werden nicht aus einer eventuell exakt null resultierenden Nettoladung
allein abgeleitet.

\section{Residualindikator für stationäre Tet4-Wärme}
Zielgröße ist die Energienorm
$\|e_T\|_E^2=\int\nabla e_T^T\mathbf k\nabla e_T$ in W K.
In einer affinen Tet4-Zelle mit konstantem Material ist
$\nabla\cdot\mathbf q_h=0$; die Quelle und inneren Flusssprünge bleiben.
Mit $\lambda_K=\lambda_{\min}(\mathbf k_K)$ und
$\lambda_F=\min(\lambda_K,\lambda_{K'})$ lautet die implementierte Gewichtung
\begin{align}\label{eq:indicator}
\eta_K^2={}&\frac{h_K^2}{\lambda_K}\int_K Q^2\,\mathrm dV
+\frac12\sum_{F\ {\rm innen}}\frac{h_F}{\lambda_F}
\int_F[\mathbf q_h\cdot\mathbf n]^2\,\mathrm dA\\
&+\sum_{F\ {\rm Neumann/Robin}}\frac{h_F}{\lambda_K}
\int_F r_F^2\,\mathrm dA,\qquad
r_F=\mathbf q_h\cdot\mathbf n-\bar q-h(T_h-T_\infty).\nonumber
\end{align}
$h_K,h_F$ sind maximale Kantenlängen. Der Sprung summiert die beiden nach
außen gerichteten Normalflüsse; innere Beiträge werden einmal berechnet und
hälftig auf beide Zellen verteilt. An Materialgrenzen vergleicht man Fluss,
nicht Gradient. Jeder Summand hat W K. Dirichlet-Flächen sind ausdrücklich
markiert und tragen keinen Randresidualterm; der vollständige äußere Rand
muss vorhanden sein. Quadraturgrade gelten für das \emph{quadrierte} Residuum.
\index{Wärmeindikator}

Die API bietet keine Zuverlässigkeits- oder Effizienzkonstanten, keine garantierte
Fehlerschranke, Datenoszillation oder Robustheit für beliebige Anisotropie und
Materialkontraste. Tet10, Transienten und Kopplungen sind ausgeschlossen.
Markierung wählt die kleinste Zahl von Zellen mit
$\sum_{K\in\mathcal M}\eta_K^2\ge\vartheta\sum_K\eta_K^2$, $0<\vartheta\le1$;
Gleichstände folgen Zellindex, eine exakte Nullsumme markiert nichts.
Eine verborgene dimensionsbehaftete Rauschschwelle oder ein Remeshing-Zyklus
wird nicht ergänzt. Beispielsweise liefern Indikatoren $(3,2,1)$ und
$\vartheta=1/2$ die erste Zelle, da $9\ge7$.

\section{Genauigkeit und Kosten sinnvoll vergleichen}
Für glatte Lösungen auf regulären Netzen werden bei P1 ungefähr
L2-Ordnung zwei und H1-Ordnung eins erwartet, bei P2 drei bzw. zwei.
Das sind asymptotische Feldaussagen, keine Garantie bei Singularitäten.
Der beobachtete Exponent ist
$p=\log(e_{h_1}/e_{h_2})/\log(h_1/h_2)$.
Eine Eigenfrequenz-/Harmonik-Zielgröße kann schneller konvergieren als die
H1-Feldnorm; die bestehenden P2-Stab-Zielgrößenordnungen um vier dürfen
nicht als Feldordnung vier ausgegeben werden.

Der unveränderte Vergleich aus Doc/54 misst Wärme und Biegung auf drei Größen,
Tet4/Tet10 und SciPy/PETSc direkt, drei Wiederholungen in frischen Prozessen.
Messzeit umfasst Solve-Pipeline, nicht Netzaufbau oder Fehlerquadratur;
RSS ist Prozesshoch, kein inkrementeller Solverbedarf. Faire Vergleiche
verwenden gleiche Fehlerziele oder explizite Zeit-/Speicherbudgets.
Nicht erreichte Ziele bleiben sichtbar. Die Rohdaten unter
\path{benchmarks/baselines/m7-37c-linux-arm64.json} zeigen ein begrenztes Raster;
es wird keine universelle Überlegenheit einer Ordnung oder eines Backends
abgeleitet und kein Benchmark neu als Leistungsmessung ausgeführt.

\subsection*{Implementierung und Nachweise.}
{\small\raggedright
\path{src/fem_toolbox/solvers/linear_backend.py},
\path{src/fem_toolbox/solvers/petsc_backend.py},
\path{src/fem_toolbox/post/heat_indicator.py};
\path{tests/unit/test_linear_backend.py},
\path{tests/unit/test_heat_indicator.py}, \path{tests/petsc},
\path{tests/convergence/test_heat_indicator_convergence.py} und
\path{tests/regression/test_accuracy_cost.py}.

\par}
